\documentclass[11pt]{article}
\usepackage[utf8]{inputenc}
\usepackage[T1]{fontenc}
\usepackage[margin=2cm]{geometry}
\usepackage{tikz}
\usepackage{tcolorbox}
\usepackage{array}
\usepackage{amsmath}
\usepackage{amssymb}
\usepackage{graphicx}
\usepackage{fancyhdr}
\usepackage{enumitem}

\tcbuselibrary{skins}
%\pagestyle{empty}

% ---- En-tête Nom / Groupe ----
\newcommand{\entete}{%
  \noindent Nom: \rule{6.5cm}{0.4pt} \hfill Groupe: \rule{3cm}{0.4pt}\\[4mm]
}

% ---- Boîte "Cible" avec titre en gras dans la boîte ----
\newtcolorbox{cible}[1]{
  colback=white, colframe=black, boxrule=0.6pt, arc=0pt,
  left=4mm, right=4mm, top=2mm, bottom=2mm,
  fonttitle=\bfseries, coltitle=black,
  before upper={\textbf{#1}\par\smallskip}
}

\begin{document}

%==================== PAGE 1 ====================
\entete
\noindent \textbf{Évaluation: Exposants et chaînes d'opérations} \hfill Chapitre 1
\vspace{2mm}
\hrule
\vspace{4mm}

\begin{cible}{Cible 1: Représenter, écrire et calculer des nombres en notation exponentielle.}

1. Trouve les puissances suivantes.
\vspace{2mm}

\begin{tabular}{@{}p{0.5cm}p{2.3cm}p{4.3cm}p{0.5cm}p{2.3cm}p{4.3cm}@{}}
a) & $2^4 = $ & \rule{3.2cm}{0.4pt} & b) & $4^3 = $ & \rule{3.2cm}{0.4pt} \\[4mm]
c) & $7^0 = $ & \rule{3.2cm}{0.4pt} & d) & $9^1 = $ & \rule{3.2cm}{0.4pt} \\[4mm]
e) & $(-4)^2 = $ & \rule{3.2cm}{0.4pt} & f) & $(-4)^3 = $ & \rule{3.2cm}{0.4pt} \\[4mm]
g) & $-(4^2) = $ & \rule{3.2cm}{0.4pt} & h) & $-3^2 = $ & \rule{3.2cm}{0.4pt} \\
\end{tabular}

\vspace{6mm}
\noindent 2. Factorise 32 en nombres premiers.

\vspace{2.5cm}

\noindent Rép: $32 = $ \rule{3.5cm}{0.4pt}


\vspace{6mm}
\noindent 3. Factorise 525 en nombres premiers.

\vspace{2.5cm}

\noindent Rép: $525 = $ \rule{3.5cm}{0.4pt}

\vspace{6mm}

\noindent 4. Factorise 351 en nombres premiers.

\vspace{2.5cm}

\noindent Rép: $351 = $ \rule{3.5cm}{0.4pt}
\vspace{4mm}
\end{cible}


\newpage

%==================== PAGE 2 ====================


\begin{cible}{Cible 2: Situer et comparer des nombres entiers.}

1. Compare les expressions suivantes à l'aide des symboles $<, >$ et $=$. 

\vspace{1cm}

\textbf{CALCULS OBLIGATOIRES}

\vspace{1cm}

\begin{tabular}{|p{7.7cm}|p{7.7cm}|}
\hline
\rule{0pt}{4mm} a) $4^3 \; \underline{\hspace{1cm}} \; 4^3$                     &               b) $(-3)^2  \; \underline{\hspace{1cm}} -3^2\; $           \\[3.2cm]
\hline
\rule{0pt}{4mm} c) $-9 \times 11 \; \underline{\hspace{1cm}} \; 3^3 \times -11$ &               d) $-20 + 40 \; \underline{\hspace{1cm}} \; -7 \times -3$ \\[3.2cm]
\hline
\end{tabular}



\vspace{6mm}
\noindent 2. Place les nombres  suivants en ordre croissant (du plus petit au plus grand): 

\vspace{2mm}
\begin{tabular}{@{}p{0.5cm}p{3cm}p{0.5cm}p{3.5cm}p{0.5cm}p{3cm}p{0.5cm}p{3cm}@{}}
a) & $(-3)^2$ & b) & $-15$ & c) & $6 + -3^2$ & d) & $(8-12)^2$ \\
\end{tabular}

\vspace{1cm} 

\noindent\rule{15cm}{0.4pt}

\vspace{6mm}
\noindent 3. Place les nombres  suivants en ordre croissant (du plus petit au plus grand): 

\vspace{2mm}
\begin{tabular}{@{}p{0.5cm}p{3cm}p{0.5cm}p{3.5cm}p{0.5cm}p{3cm}p{0.5cm}p{3cm}@{}}
a) & $0^1$ & b) & $1^0$ & c) & $-2^1$ & d) & $(-2)^2$ \\
\end{tabular}

\vspace{1cm}

\noindent\rule{15cm}{0.4pt}

\end{cible}

\vspace{5mm}



\newpage

%==================== PAGE 3 ====================


\begin{cible}{Cible 3: Effectuer des chaînes d'opérations.}

1. Trouve le résultat de la chaine d'opérations suivante. Écris tes étapes de calcul.

\vspace{1cm}

\renewcommand{\arraystretch}{1}
\begin{tabular}{|p{7.7cm}|p{7.7cm}|}
\hline
\rule{0pt}{4mm} a) $4 \times (6+3) = $ & b) $6 - (-7)^2 = $ \\[8cm]
\hline
\rule{0pt}{4mm} c) $ (-3)^3 - 7 = $ & d) $5 \times -4 + 9 = $ \\[8cm]
\hline
\end{tabular}

\end{cible}

\newpage

\begin{cible}{Cible 3: Effectuer des chaînes d'opérations.(suite)}

1. Trouve le résultat de la chaine d'opérations suivante. Écris tes étapes de calcul.

\vspace{1cm}

\renewcommand{\arraystretch}{1}
\begin{tabular}{|p{7.7cm}|p{7.7cm}|}
\hline
\rule{0pt}{4mm} e) $5 \times 9 - 50 \div (12-7) = $ & f) $ 4 \times 6 - (-2^3 = $ \\[8cm]
\hline
\rule{0pt}{4mm} g) $(-7+5)^3 \div -4 = $ & h) $(-7)^2 + 8 - 4 \div 4 = $ \\[8cm]
\hline
\end{tabular}

\end{cible}


\end{document}